\subsection{Local input design improves declared information criteria}
All parameter-design cases completed after the shutdown-grid correction, and
the independent audit passed the information matrices, eigenpairs, objective
values, selected-start rule and resource constraints. Each selected local
solution improves its own declared reference criterion. This comparison is
conditional on the fitted model points, working covariance and initialization;
it does not select a physically correct model.

The ensemble designs show how the target changes the proposed experiment
(Table~\ref{tab:parameter-design} and Fig.~\ref{fig:design-utilities}). The
minimum criterion targets the least favorable relative scenario; the mean
criterion targets average log improvement. Historical information changes both the objective
and the reference: future-only ensembles use multisine because constant input
leaves delay unobservable in equilibrated scenarios, whereas adding the two
historical calibration records makes the constant-reference total information
nonsingular. Differences between the two panels therefore cannot be read as
absolute information loss from including historical data. The determinant,
minimum-eigenvalue, delay and prediction criteria also have different scales.
% Generated by plot_design_results.py from separately audited phase summaries.
\begin{table}[tbp]
\centering\small
\begin{tabular}{llrr}
\toprule
Ensemble, information & Criterion & $\Delta\log$ utility & Max. (${}^\circ$C) \\
\midrule
mean, future & D & 1.508 & 84.990 \\
mean, history & D & 1.390 & 84.990 \\
mean, future & E & 1.423 & 74.333 \\
mean, history & E & 0.754 & 80.792 \\
mean, future & Delay & 1.439 & 74.333 \\
mean, history & Delay & 0.758 & 80.792 \\
mean, future & Prediction & 0.861 & 84.990 \\
mean, history & Prediction & 0.348 & 84.990 \\
min, future & D & 1.294 & 84.990 \\
min, history & D & 1.299 & 84.151 \\
min, future & E & 0.560 & 84.990 \\
min, history & E & 0.621 & 75.349 \\
min, future & Delay & 0.620 & 73.829 \\
min, history & Delay & 0.631 & 75.463 \\
min, future & Prediction & 0.704 & 84.990 \\
min, history & Prediction & 0.300 & 84.990 \\
\bottomrule
\end{tabular}
\caption{Selected local ensemble designs over 48 model/startup/covariance scenarios. Utility is the mean or minimum scenario-relative log improvement. Future-only ensembles use multisine as reference; ensembles with historical information use constant input, whose total information is nonsingular after adding calibration information. Criterion scales differ. Temperature is the largest independently checked value across sixteen fitted mean models through the shutdown tail. These local results do not establish global optimality, probabilistic model robustness or hardware limits.}
\label{tab:parameter-design}
\end{table}


Figure~\ref{fig:design-commands} displays two selected ensemble determinant
designs and their dense temperature checks. The upper temperature envelope
ranges over the declared sixteen fitted mean models, without probabilities
over models or parameter uncertainty. The selected inputs use the amplitude
and slew limits and approach the imposed temperature constraint. The finite
tail check includes the last command's ramp to zero. All selected designs pass
the independent dense acceptance check, but this does not establish a
continuous-time or hardware temperature bound.

\begin{figure}[p]
 \centering
 \includegraphics[width=\linewidth]{design_commands_temperature.pdf}
 \caption{Information design changes the command and its predicted thermal
 envelope. Constant and multisine are common candidates; ensemble mean D and
 min D are selected determinant designs using future information only. Commands
 share the declared horizon, time-average dose, amplitude and slew constraints.
 The lower panel is the pointwise maximum temperature among sixteen fitted
 mean/startup combinations, not a confidence band or one physical trajectory.
 The vertical line marks the end of the design horizon and start of shutdown;
 the horizontal line is the assumed temperature acceptance limit. Calculations
 use a $0.25$ s audit grid through the $300$ s tail.}
 \label{fig:design-commands}
\end{figure}

\begin{figure}[p]
 \centering
 \includegraphics[width=\linewidth]{design_relative_utilities.pdf}
 \caption{The design target and scenario aggregate change relative utility.
 Points are the independently audited selected local objective values. Mean
 and min aggregate scenario-relative log improvements over 48 fixed
 model/startup/covariance scenarios. The first panel uses multisine as reference;
 the second adds historical calibration information and uses constant input.
 Criteria have different targets and scales, so ordinate heights across
 criteria or reference changes do not rank absolute physical information.
 Lines join categorical criteria as guides. No global optimality certificate
 is implied.}
 \label{fig:design-utilities}
\end{figure}

\subsection{Refitting competitors changes model discrimination}
The separate fixed-mean discrimination designs compare all six pairs among
the representative formulations (Table~\ref{tab:discrimination-design}). Their
mean/min objectives concern log improvements relative to multisine for each
pair. The smallest raw pair distance is a different summary and need not rank
those objectives in the same order. These designs share the parameter-design
resource assumptions but optimize mean separation rather than parameter
information.
% Generated by plot_design_results.py from separately audited phase summaries.
\begin{table}[tbp]
\centering\small
\begin{tabular}{llrr}
\toprule
Startup, covariance & Aggregate & $\Delta\log D^2$ & Smallest $D^2$ \\
\midrule
Cold, IID & mean & 0.852 & 52.329 \\
Cold, IID & min & 0.798 & 52.329 \\
Cold, OU10 & mean & 0.825 & 17.797 \\
Cold, OU10 & min & 0.652 & 18.054 \\
Cold, OU60 & mean & 0.949 & 59.994 \\
Cold, OU60 & min & 0.628 & 58.797 \\
Equilibrated, IID & mean & 1.184 & 18.955 \\
Equilibrated, IID & min & 1.246 & 19.525 \\
Equilibrated, OU10 & mean & 1.451 & 2.136 \\
Equilibrated, OU10 & min & 1.324 & 1.564 \\
Equilibrated, OU60 & mean & 2.596 & 6.202 \\
Equilibrated, OU60 & min & 2.368 & 6.202 \\
\bottomrule
\end{tabular}
\caption{Fixed-mean discrimination designs compare six pairs among the four representative formulations. The objective aggregates pairwise log improvements relative to multisine, whereas the last column is the smallest raw selected pair distance; these are different comparisons. Covariances are common known working laws with marginal standard deviation $0.25\,{}^\circ\mathrm C$. Alternative-parameter refitting is evaluated separately in Fig.~\ref{fig:discrimination-refits}.}
\label{tab:discrimination-design}
\end{table}


Allowing the competing formulation to change parameters reduces its distance
from the fixed one-stage generating mean
(Fig.~\ref{fig:discrimination-refits}). Under equilibrated constant input, the
closest-found distance vanishes to numerical precision for every alternative
and working covariance. The matched gain and time constant give an analytic
zero-distance witness because the delay factor is absent. Positive separation
at the original fitted alternative parameters had therefore overstated the
information in that input about model order.

The selected varying inputs retain larger closest-found distances than
multisine in this bounded directional comparison, including after refitting.
That result concerns a fixed one-stage generating mean against three
alternatives, rather than every possible generating formulation or parameter
point. Each smallest converged distance is an upper bound on the global closest
alternative. The ideal equal-prior Gaussian error formula for two fixed means
cannot be reused as a power calculation after fitting those alternatives or
allowing the observed between-record discrepancy.

\begin{figure}[p]
 \centering
 \includegraphics[width=\linewidth]{discrimination_refits.pdf}
 \caption{Parameter refitting can remove apparent model-order separation.
 The fixed one-stage generating mean is compared with the three indicated
 alternative formulations under the first declared covariance, IID Gaussian
 error with marginal standard deviation $0.25\,{}^\circ\mathrm C$. Rows use
 cold and input-equilibrated delay states. Left panels retain each alternative's
 original fitted parameters; right panels use its smallest converged
 24-start distance. Const. and Multi. are common candidates; Mean and Min are
 the selected fixed-mean discrimination inputs. Symmetric-log axes are linear
 below one, retaining vanishing distances without a positive plotting floor.
 Lines join categorical waveforms. The complete saved comparison also includes
 both correlated working covariances. Distances after refitting are upper
 bounds on global minima, without a real-data testing-power interpretation.}
 \label{fig:discrimination-refits}
\end{figure}
